\documentclass[10pt]{article}

\usepackage[T1]{fontenc}
\usepackage[margin=0.68in]{geometry}
\usepackage{graphicx}
\usepackage{booktabs}
\usepackage{tabularx}
\usepackage{array}
\usepackage{xcolor}
\usepackage{hyperref}
\usepackage{microtype}
\usepackage{enumitem}
\usepackage[font=small,labelfont=bf]{caption}

\definecolor{accent}{HTML}{1F4E79}
\definecolor{lightgray}{HTML}{F2F4F7}
\hypersetup{colorlinks=true,linkcolor=accent,citecolor=accent,urlcolor=accent}
\setlength{\parindent}{0pt}
\setlength{\parskip}{3.5pt}
\setlist{nosep,leftmargin=1.25em}
\renewcommand{\arraystretch}{1.08}

\title{\vspace{-2.0em}\textbf{CSE 8803 HW1: Building a Decoder-Only Transformer from Scratch}}
\author{GT Username: \texttt{yzhu751}}
\date{Georgia Institute of Technology -- Fall 2026}

\begin{document}
\maketitle
\vspace{-1.2em}

\begin{abstract}
For this project, I built a byte-level BPE tokenizer and a decoder-only language model without
using ready-made tokenizer, attention, or Transformer modules. The model includes RMSNorm, rotary
position embeddings (RoPE), grouped-query attention (GQA), SwiGLU, tied embeddings, my own AdamW
optimizer, and a KV cache for faster step-by-step generation. I trained the required 29.37M-
parameter GQA model for 3,000 steps on TinyStories. On the validation data, it reached a perplexity
of 5.403 and 0.497 bits per byte (BPB). I checked the implementation with public tests, staff
reference files, exact tokenizer and data hashes, cached and uncached comparisons, and a small
overfitting test.
\end{abstract}

\section{Implementation and experimental setup}

The tokenizer starts with all 256 possible byte values and one special
\texttt{<|endoftext|>} token. It then learns 7,935 fixed pair merges, giving a final vocabulary of
8,192 tokens. Its SHA-256 hash exactly matched the official checksum, which shows that the learned
tokenizer was correct. The language model has 8 decoder layers, with normalization before each
main block. It uses $d_{\mathrm{model}}=512$, 8 query heads, head size 64, a SwiGLU hidden size of
1,536, and a context length of 512. The required GQA model uses 4 KV heads. For a fair comparison,
the MHA model changes only the number of KV heads from 4 to 8. RoPE is added to queries and keys,
and cached decoding uses the correct position offset \cite{su2021roformer}. The RMSNorm, SwiGLU,
and GQA designs follow \cite{zhang2019rmsnorm,shazeer2020glu,ainslie2023gqa}.

I used random seed 2026 and BF16 number format. Each small batch contained 8 sequences, and I
combined gradients across 32 small batches. Since every sequence had 512 tokens, each optimizer
step used $8\times32\times512=2^{17}$ tokens. AdamW used $\beta=(0.9,0.95)$, weight decay 0.1,
and a maximum gradient norm of 1.0 before clipping. The learning rate rose in a straight line for
150 steps until it reached $3\times10^{-4}$, then followed a cosine curve down to
$3\times10^{-5}$. The GQA run saw 393,216,000 tokens during 3,000 steps on one NVIDIA H200 with
PyTorch 2.13.0+cu126. Training took 2,011.7 seconds (0.559 h), averaged 195,467 tokens/s, and used
at most 2.65 GiB of allocated GPU memory. The MHA model trained for 1,000 steps on an NVIDIA H100.
Because the two training runs used different GPU models, I do not use their training speeds as a
fair GQA-versus-MHA comparison. I ran the cache tests for both checkpoints one after the other on
the same H200.

\section{Tokenizer analysis}

Table~\ref{tab:tokenizer} tests shorter and longer prefixes of the same learned merge list. I did
not train a new tokenizer for every vocabulary size. A larger vocabulary compressed the text
better: bytes per token almost doubled from 2.480 to 4.895, while tokens per word dropped from
2.040 to 1.034. The tradeoff was speed. Encoding slowed from 0.753 to 0.546 MB/s because the
tokenizer had more merge rules to check. Learning the full vocabulary from the fixed 46,906-byte
data shard took 0.234 s on ICE, which is below the published limit of 0.853 s.

\begin{table}[ht]
\centering
\small
\caption{Tokenizer prefix analysis on the fixed 44,843-byte validation span.}
\label{tab:tokenizer}
\begin{tabular}{rrrrr}
\toprule
Vocab & Tokens & Bytes/token & Tokens/word & Encoding MB/s \\
\midrule
512  & 18,080 & 2.480 & 2.040 & 0.753 \\
1,024 & 13,234 & 3.388 & 1.494 & 0.630 \\
2,048 & 10,710 & 4.187 & 1.209 & 0.577 \\
4,096 & 9,477  & 4.732 & 1.070 & 0.552 \\
8,192 & 9,161  & 4.895 & 1.034 & 0.546 \\
\bottomrule
\end{tabular}
\end{table}

\begin{figure}[ht]
\centering
\includegraphics[width=0.62\linewidth]{../tmp/pdfs/tokenizer_vocabulary.pdf}
\caption{Vocabulary size compared with bytes per token. The extra compression becomes smaller
after 4,096 tokens.}
\label{fig:tokenizer}
\end{figure}

\section{GQA training and final evaluation}

Figure~\ref{fig:training} shows that training stayed stable. Training loss and validation NLL stay
close together, so there is no clear sign of overfitting near the end. After the noisy beginning,
the gradient norm before clipping settles near 0.53. Speed stays around 190--202k tokens/s most of
the time. The short drop near step 900 did not happen with any unusual change in loss or gradient
size, so it was most likely caused by normal changes in the computer system rather than the model.

\begin{figure}[ht]
\centering
\includegraphics[width=0.96\linewidth]{../runs/gqa/training_curves.pdf}
\caption{GQA training results. The light lines show raw loss and gradient norm, while the dark
lines show 50-step moving averages. Validation NLL was measured every 100 steps.}
\label{fig:training}
\end{figure}

I evaluated the final checkpoint on a fixed piece of validation text containing 9,161 tokens, or
44,843 bytes. Each of the 9,160 next-token answers was scored exactly once. The results are shown
in Table~\ref{tab:finaleval}. Token perplexity is only fair to compare when the tokenizer stays the
same, because different token boundaries change both the number and difficulty of the predictions.
Bits per byte (BPB) is more useful across tokenizers because it divides the total negative
log-likelihood by the number of original bytes:
\[
\mathrm{BPB}=\frac{\sum_i -\log p(x_i\mid x_{<i})}{N_{\mathrm{bytes}}\log 2},
\]
This makes BPB a more meaningful way to compare different tokenizations of the same text.

\begin{table}[ht]
\centering
\small
\caption{Final GQA evaluation at step 3,000. Lower is better for all three loss metrics.}
\label{tab:finaleval}
\begin{tabular}{rrrr}
\toprule
Mean NLL & Token PPL & Bits/byte & Target tokens \\
\midrule
1.68699 & 5.40318 & 0.49715 & 9,160 \\
\bottomrule
\end{tabular}
\end{table}

For text generation, I used the four provided prompts and generated 128 new tokens for each one.
The settings were temperature 0.8, top-$k=50$, top-$p=0.95$, and seeds 2026--2029. The complete,
unchanged outputs are in Appendix~\ref{app:samples}. Each sample starts with readable
TinyStories-style writing, but two problems appear more than once. \textbf{(1) The model does not
always stop or stay on topic.} Three samples produce \texttt{<|endoftext|>} and then begin a new,
unrelated story because generation continues until it reaches the fixed token limit. Some samples
also stop in the middle of a sentence. \textbf{(2) The model loses track of characters and causes.}
For example, the robot is called ``the robot,'' ``it,'' and ``he.'' In the trumpet story, ``girl''
later changes to ``his.'' The box story forgets about the stuck box and suddenly moves to opening a
house door and finding toys. These mistakes suggest that the model can write short, grammatical
parts but has trouble remembering the full story over a longer distance.

\section{MHA versus GQA and KV-cache behavior}

In the controlled comparison after 1,000 steps, MHA has 2,097,152 more parameters, or 7.14\% more
than GQA. It also has slightly lower validation NLL and perplexity, as shown in
Table~\ref{tab:quality}. In other words, the extra KV projection weights gave MHA a small quality
advantage when both models saw the same number of training tokens. GQA still reached almost the
same quality while using fewer parameters.

\begin{table}[ht]
\centering
\small
\caption{Controlled validation comparison at step 1,000 (131,072,000 tokens seen).}
\label{tab:quality}
\begin{tabular}{lrrr}
\toprule
Model & Parameters & Validation NLL & Validation PPL \\
\midrule
GQA (8 Q, 4 KV) & 29,368,832 & 2.01174 & 7.47628 \\
MHA (8 Q, 8 KV) & 31,465,984 & 1.99206 & 7.33059 \\
\bottomrule
\end{tabular}
\end{table}

The KV cache saves old keys and values so the model does not need to compute them again for every
new token. If there are $L$ layers, batch size $B$, $T$ cached tokens, $H_{kv}$ KV heads, head size
$d_h$, and $s$ bytes per number, storing both keys and values needs
\[
M_{KV}=2LBTH_{kv}d_hs.
\]
MHA gives every query head its own key and value heads, so $H_{kv}=H_q$ and
$M_{KV}^{\mathrm{MHA}}=2LBTd_{\mathrm{model}}s$. GQA shares key and value heads across groups. If
there are $G$ groups, its cache uses only $G/H_q$ as much memory as MHA. MQA goes further by using
one KV head, so it uses $1/H_q$ as much. In this experiment, $L=8$, $T=256$, $d_h=64$, and each
BF16 number uses $s=2$ bytes. Therefore, 4-head GQA uses exactly half as much cache memory as
8-head MHA.

\begin{table}[ht]
\centering
\small
\caption{H200 greedy-decoding benchmark (128-token prompt plus 128 generated tokens).}
\label{tab:cache}
\begin{tabular}{llrrr}
\toprule
Model & Batch & Cache & Cached tokens/s & Speedup \\
\midrule
GQA & 1  & 2 MiB   & 198.42  & 0.996$\times$ \\
GQA & 32 & 64 MiB  & 6,230.97 & 1.020$\times$ \\
MHA & 1  & 4 MiB   & 210.04  & 0.995$\times$ \\
MHA & 32 & 128 MiB & 6,631.57 & 1.032$\times$ \\
\bottomrule
\end{tabular}
\end{table}

The measured speed difference is small. With a model this size, cached generation starts many tiny
GPU jobs, one token at a time. Uncached generation repeats more math, but it does that work in
larger matrix operations that use the GPU well. On an H200, the cost of starting many small jobs can
cancel out some of the cache's saved math. The clearest result is memory use: GQA cuts KV-cache
storage in half. I checked correctness with staff reference logits, identical cached and uncached
tokens, correct RoPE positions, and tests showing that after the prompt, only new tokens are sent
through the key and value projections.

\section{Required conceptual analysis and debugging evidence}

\textbf{Causal masking and parallel training.} During teacher-forced training, the whole correct
input sequence is already available. This means the model can calculate Q, K, and V for every
position at the same time. However, a token must not look ahead at future answers. A triangular
mask solves this problem by setting every future attention score ($j>i$) to $-\infty$ before
softmax. Those future positions then get probability zero. All rows are still calculated together,
but position $i$ can only use positions up to $i$, so future information cannot leak into training.

\textbf{AdamW versus L2 under Adam.} Standard L2 regularization adds $\lambda\theta$ to the
gradient before Adam calculates its moving averages and per-parameter scaling. As a result, the
amount of shrinking can change based on Adam's internal statistics. AdamW keeps weight decay
separate and directly applies $\theta\leftarrow(1-\eta\lambda)\theta$
\cite{loshchilov2019adamw}. I used parameter groups so matrix weights receive weight decay, while
the one-dimensional RMSNorm scale parameters do not.

\textbf{A real bug I fixed.} At first, the test for pausing and restarting the optimizer failed,
even though the saved state loaded without an error. By checking where the tensors were stored in
memory, I found that PyTorch's \texttt{load\_state\_dict} left the loaded
\texttt{exp\_avg} and \texttt{exp\_avg\_sq} tensors sharing storage with the original optimizer
state. When the original optimizer took another step, it secretly changed the state that the
restarted optimizer was about to use. This made their next parameters different. Loading a
\texttt{deepcopy(state\_dict)} gave the restarted optimizer its own separate tensors and fixed the
problem. After that, both the parameter-group restart test and the full uninterrupted-versus-
restarted training test passed. The second test also confirmed that the random-number and data-
generator states were restored correctly.

\section{Conclusion}

This project completes the required tokenizer and language-model training process from beginning
to end. The final GQA model is close to the published reference quality while using only half as
much KV-cache memory as MHA. The experiments also show that doing less math does not always make a
program much faster: for a small model generating one token at a time on a powerful GPU, the extra
overhead of many small GPU jobs can hide the speed benefit.

\clearpage
\appendix
\section{Full generated samples}
\label{app:samples}
\small

\subsection{Prompt: Once upon a time (seed 2026)}
Once upon a time, a little girl named Lucy went to the garden. She saw her friend Tom. Tom was
eating some broccoli. Lucy was jealous. She wanted a big broccoli and wanted a small broccoli.
Lucy said, ``Please, can we both have a small piece of broccoli?'' Tom nodded and said yes. They
shared the broccoli and they were happy. From that day on, Lucy and Tom shared the broccoli with
each other. \texttt{\textless|endoftext|\textgreater} One day, a little girl named Lily went to the
park with her mom. They saw a big tree with a long string hanging from it. Lily wanted to get the
string, so she ran to the tree and got her long necklace

\subsection{Prompt: Lily found a strange box (seed 2027)}
Lily found a strange box and she wondered what was inside. She tried to open it, but it was stuck.
She tried and tried to open it, but it wouldn't budge. She felt sad and frustrated. Suddenly, she
had an idea. She ran to the house and opened the door. She opened it and saw many toys inside. They
were not normal to play with anymore. They were not so bad after all. She was so happy, she hugged
the toy and played with it all day. \texttt{\textless|endoftext|\textgreater} Sam and Ben were
playing with their toys. Sam had a big truck and Ben had a small car. They liked to make the car go
fast on the floor and pretend it was their car. Sam had an idea. He thought it was a new game. He
wanted to make it go

\subsection{Prompt: The little robot looked at the moon (seed 2028)}
The little robot looked at the moon and said, ``This place is too dark. You should go back home.''
The robot was very sad. It wished he had stayed home with its mommy and daddy. Sadly, it was too
late. The robot was gone forever. \texttt{\textless|endoftext|\textgreater} Once upon a time, there
was a girl who wanted to learn how to play the trumpet. She went to her room and grabbed her
trumpet. She started to play a song to it. She thought it was funny and made her feel strange.
Suddenly, a gust of wind blew and the trumpet went flying out of his room. It was blocking the whole
room! The girl was very sad and she

\subsection{Prompt: Ben and Mia went into the forest (seed 2029)}
Ben and Mia went into the forest. They saw many trees and flowers. They saw a fountain and a
squirrel. They saw a pond with ducks and a squirrel on them. They saw a rabbit who was sleeping on
a tree. They saw a hole in the ground and ran towards it. They wanted to hide in it, but they heard
a voice. ``Hey, you two, get out of there! It's my home! It's my home!'' It was a small bird who had
seen them from the hole in the fence. The bird looked sad and hungry. She saw Ben and Mia and felt
sorry for them. She decided to be gentle and

\clearpage
\begin{thebibliography}{9}
\bibitem{sennrich2016bpe}
R. Sennrich, B. Haddow, and A. Birch. Neural Machine Translation of Rare Words with Subword Units.
ACL, 2016.

\bibitem{vaswani2017attention}
A. Vaswani et al. Attention Is All You Need. NeurIPS, 2017.

\bibitem{su2021roformer}
J. Su et al. RoFormer: Enhanced Transformer with Rotary Position Embedding. 2021.

\bibitem{zhang2019rmsnorm}
B. Zhang and R. Sennrich. Root Mean Square Layer Normalization. NeurIPS, 2019.

\bibitem{shazeer2020glu}
N. Shazeer. GLU Variants Improve Transformer. 2020.

\bibitem{ainslie2023gqa}
J. Ainslie et al. GQA: Training Generalized Multi-Query Transformer Models from Multi-Head
Checkpoints. EMNLP, 2023.

\bibitem{loshchilov2019adamw}
I. Loshchilov and F. Hutter. Decoupled Weight Decay Regularization. ICLR, 2019.
\end{thebibliography}

\end{document}
